Deployed large language model (LLM) systems face continuously evolving input distributions, yet standard drift detectors rely on low-order statistics that may miss structural reorganization of the embedding manifold. We investigate whether persistent-homology-based drift detectors can complement classical methods on real LLM embedding streams. In a controlled comparative study across two datasets (\agnews, \twentyng), two embedding models (\minilm 384-d, \bertbase 768-d), and 11 drift scenarios---including novel centroid-preserving drifts (within-topic subtopic reweighting and style perturbation)---we evaluate 20+ detection methods spanning statistical baselines (centroid shift, covariance shift, MMD, energy distance, classifier-based two-sample test), TDA scalar features, persistence landscapes, and sliced Wasserstein distances on persistence diagrams. On standard topic drift, classical methods achieve near-perfect AUC. On harder centroid-preserving scenarios, TDA methods provide complementary signal. In a controlled synthetic experiment evaluated with proper AUC (replacing separability ratios), H1 persistent entropy uniquely detects loop emergence where centroid shift fails entirely. Ablation studies over window size (50--400), TDA subsample size (40--160), and PCA dimensionality (20--384) reveal that PCA-50 projection before persistent homology substantially improves TDA reliability in high-dimensional embedding spaces. We recommend a two-tier monitoring architecture combining fast statistical primary detectors with TDA secondary detectors for topology-sensitive drift at modest computational overhead.
